% answers_short.tex —— 2002 年·北京卷·理科 · 简短答案速查
% 由 tools/make_short_answers.py 自 answers.tex 的 \gksolhead/\gkanswer 提取，
% 并与 2002/2002_数学真题_北京理_答案要点.md 逐题比对后生成，勿手改（改源 MD 后重跑）。

\begin{gknotice}
\textbf{使用说明}：本节只给最终答案便于快速核对，完整推导见解析版。答案为编写者独立推导，\textbf{非官方标准答案}。
\end{gknotice}

\gkpart{选择题}
\begin{center}\begin{tabular}{ccccccccccccc}
\toprule
题号 & 1 & 2 & 3 & 4 & 5 & 6 & 7 & 8 & 9 & 10 & 11 & 12 \\
\midrule
答案 & B & D & B & C & D & D & C & A & A & C & B & A \\
\bottomrule
\end{tabular}\end{center}

\gkpart{填空题}
\begin{itemize}[leftmargin=2.6em]
  \item[\textbf{13.}] \(\arctan(-\dfrac54)<\arcsin(-\dfrac25)<\arccos(-\dfrac34)\)。
  \item[\textbf{14.}] 4。
  \item[\textbf{15.}] ①②③④⑤。
  \item[\textbf{16.}] U/PENDING。
\end{itemize}

\gkpart{解答题}
\begin{itemize}[leftmargin=2.6em]
  \item[\textbf{17.}] \([\dfrac12,5)\)。
  \item[\textbf{18.}] (1) \(\theta=\arctan\dfrac{2h}{b-d}\)；(2) \(EF\parallel ABCD\)；(3) \(V_{\text{估}}<V\)
  \item[\textbf{19.}] (1) \(x_n\geqslant\sqrt a\quad(n\geqslant2)\)；(2) \(x_n\geqslant x_{n+1}\quad(n\geqslant2)\)；(3) \(\displaystyle\lim_{n\to\infty}x_n=\sqrt a\)
  \item[\textbf{20.}] (1) 至少2个单位时间；一种填表方法见解析；(2) 7个单位时间
  \item[\textbf{21.}] (1) \(G(\frac{1+b}{3},\frac c3)\)，\(F(\frac12,\frac{b^2+c^2-b}{2c})\)，\(H(b,\frac{b(1-b)}c)\)；三点共线；(2) \(c^2=3b(1-b)\)，\(0<b<1\)，并排除 \((b,c)=(\frac12,\pm\frac{\sqrt3}{2})\)
  \item[\textbf{22.}] (1) \(f(0)=f(1)=0\)；(2) \(f(x)\) 为奇函数；(3) \(\displaystyle S_n=2^{-n}-1\)
\end{itemize}
